\RequirePackage[l2tabu, orthodox]{nag}
\RequirePackage{silence}
\documentclass[french,english]{beamer}
%INSTALL

%avoids a warning
\usepackage[log-declarations=false]{xparse}
\usepackage{fontspec} %font selecting commands
\usepackage{xunicode}
\usepackage{ucharclasses}
%warn about missing characters
\tracinglostchars=2

%REDAC
\usepackage{booktabs}
\usepackage{calc}
\usepackage{tabularx}

\usepackage{mathtools} %load this before babel!
	\mathtoolsset{showonlyrefs,showmanualtags}

\usepackage[french, english]{babel}
%suppresses the warning about frenchb not modifying the captions (“—” to “:” in “Figure 1 – Legend”).
	\frenchbsetup{AutoSpacePunctuation=false,SuppressWarning=true}

%\usepackage[super]{nth}%better use fmtcount! (loaded by datetime anyway; see below about pbl with warnings and package silence)
\usepackage{listings} %typeset source code listings
	\lstset{language=XML,tabsize=2,captionpos=b,basicstyle=\NoAutoSpacing}%NoAutoSpacing avoids space before colon or ?}%,literate={"}{{\tt"}}1, keywordstyle=\fontspec{Latin Modern Mono Light}\textbf, emph={String, PreparedStatement}, emphstyle=\fontspec{Latin Modern Mono Light}\textbf, language=Java, basicstyle=\small\NoAutoSpacing\ttfamily, aboveskip=0pt, belowskip=0pt, showstringspaces=false
\usepackage[nolist,smaller,printonlyused]{acronym}%,smaller option produces warnings from relsize in some cases, it seems.% Note silence and acronym and hyperref make (xe)latex crash when ac used in section (http://tex.stackexchange.com/questions/103483/strange-packages-interaction-acronyms-silence-hyperref), rather use \section{\texorpdfstring{\acs{UE}}{UE}}.
\usepackage{fmtcount}
\usepackage[nodayofweek]{datetime}%must be loaded after the babel package. However, loading it after {nth} generates a warning from fmtcount about ordinal being already defined. Better load it before nth? (then we can remove the silence package which creates possible crashes, see above.) Or remove nth?
%\usepackage{xspace}%do we need this?
\usepackage{hyperref}% option pdfusetitle must be introduced here, not in hypersetup.
%breaklinks makes links on multiple lines into different PDF links to the same target.
%colorlinks (false): Colors the text of links and anchors. The colors chosen depend on the the type of link. In spite of colored boxes, the colored text remains when printing.
%linkcolor=black: this leaves other links in colors, e.g. refs in green, don't print well.
%pdfborder (0 0 1, set to 0 0 0 if colorlinks): width of PDF link border
%hidelinks
\hypersetup{breaklinks, bookmarksopen, hyperfigures=true}
%in Beamer, sets url colored links but does not change the rest of the colors (http://tex.stackexchange.com/questions/13423/how-to-change-the-color-of-href-links-for-real)
%\hypersetup{breaklinks,bookmarksopen,colorlinks=true,urlcolor=blue,linkcolor=,hyperfigures=true}
% hyperref doc says: Package bookmark replaces hyperref’s bookmark organization by a new algorithm (...) Therefore I recommend using this package.
\usepackage{bookmark}

% center floats by default, but do not use with float
% \usepackage{floatrow}
% \makeatletter
% \g@addto@macro\@floatboxreset\centering
% \makeatother
\usepackage{ragged2e} %new com­mands \Cen­ter­ing, \RaggedLeft, and \RaggedRight and new en­vi­ron­ments Cen­ter, FlushLeft, and FlushRight, which set ragged text and are eas­ily con­fig­urable to al­low hy­phen­ation (the cor­re­spond­ing com­mands in LaTeX, all of whose names are lower-case, pre­vent hy­phen­ation al­to­gether). 
\usepackage{siunitx} %[expproduct=tighttimes, decimalsymbol=comma] ou (plus récent ?) [round-mode=figures, round-precision=2, scientific-notation = engineering]
\sisetup{detect-all, locale = FR, strict}% to detect e.g. when in math mode (use a math font) - check whether this makes sense with strict
\usepackage{braket} %for \Set
\usepackage{natbib}
\usepackage{doi}

\usepackage{amsmath,amsthm}
% \usepackage{amsfonts} %not required?
% \usepackage{dsfont} %for what?
%unicode-math overwrites the following commands from the mathtools package: \dblcolon, \coloneqq, \Coloneqq, \eqqcolon. Using the other colon-like commands from mathtools will lead to inconsistencies. Plus, Using \overbracket and \underbracke from mathtools package. Use \Uoverbracket and \Uunderbracke for original unicode-math definition.
%use exclusively \mathbf and choose math bold style below.
\usepackage[warnings-off={mathtools-colon, mathtools-overbracket}, bold-style=ISO]{unicode-math}

\defaultfontfeatures{
	Fractions=On,
	Mapping=tex-text% to turn "--" into dashes, useful for bibtex%%
}
%\defaultfontfeatures[\rmfamily]{
%	Fractions=On,
%	Mapping=% to leave " alone (disable the default mapping tex-text; requires loading the font afterwards?)
%}
\newfontfamily\xitsfamily{XITS}
\newfontfamily\texgyretermesfamily{TeX Gyre Termes}
\newfontfamily\lmfamily{Latin Modern Roman}
\setmainfont{Latin Modern Roman}
\setsansfont{Latin Modern Sans}
\setmonofont{Latin Modern Mono}
\defaultfontfeatures{}%disable default font features to avoid warnings with math fonts.
\setmathfont{XITS Math}
\setmathfont[range={\mathcal,\mathbfcal},StylisticSet=1]{XITS Math}
%silently delays restoring normal font when using e.g. “≠)” or “≠ other word”.
%\setTransitionTo{Arrows}{\xitsfamily}
%\setTransitionTo{MathematicalOperators}{\xitsfamily}
%this solution explicitly fails on “≠)” and works for “≠ other word.”
\setTransitionsFor{MathematicalOperators}{\begingroup\xitsfamily}{\endgroup}
\setTransitionsFor{Arrows}{\begingroup\xitsfamily}{\endgroup}
%texgyretermes does not have ✓, XITS does.
\setTransitionsFor{Dingbats}{\begingroup\xitsfamily}{\endgroup}
%this also works, but it is more complex
%\def\ResetTransitionTo#1{%
%  \XeTeXinterchartoks 255 \csname#1Class\endcsname{\relax}}
%\setTransitionsFor{MathematicalOperators}
%  {\begingroup\ResetTransitionTo{MathematicalOperators}\xitsfamily}
%  {\endgroup}

\usepackage{cleveref}% cleveref should go "laster" than hyperref
%GRAPHICS
\usepackage{pgf}
\usepackage{pgfplots}
	\usetikzlibrary{babel, matrix, fit, plotmarks, calc, trees, shapes.geometric, positioning, plothandlers, arrows, shapes.multipart}
\pgfplotsset{compat=1.11}
\usepackage{graphicx}

\graphicspath{{graphics/},{graphics-dm/}}
\DeclareGraphicsExtensions{.pdf}
\LetLtxMacro\SavedIncludeGraphics\includegraphics
\AtBeginDocument{
	\def\includegraphics#1#{% #1 catches optional stuff (star/opt. arg.)
		\IncludeGraphicsAux{#1}%
	}%
}
\newcommand*{\IncludeGraphicsAux}[2]{%
	\XeTeXLinkBox{%
		\SavedIncludeGraphics#1{#2}%
	}%
}%

%HACKING
\usepackage{printlen}
\uselengthunit{mm}
% 	\newlength{\templ}% or LenTemp?
% 	\setlength{\templ}{6 pt}
% 	\printlength{\templ}
\usepackage{etoolbox} %for addtocmd command
\usepackage{scrhack}% load at end. Corrects a bug in float package, which is outdated but might be used by other packages
\usepackage{xltxtra} %somebody said that this is loaded by fontspec, but does not seem correct: if not loaded explicitly, does not appear in the log and \showhyphens is not corrected.

\usepackage{appendixnumberbeamer}
\setbeamertemplate{navigation symbols}{} 
\usepackage{preamble/beamerthemeParisFrance}
\usefonttheme{professionalfonts}
\setcounter{tocdepth}{10}
%From: http://tex.stackexchange.com/questions/168057/beamer-with-xelatex-on-texlive2013-enumerate-numbers-in-black
%I don’t think it’s useful to submit this as a bug: nothing has been solved since March, 2015. See: https://bitbucket.org/rivanvx/beamer/issues?status=resolved.

\setbeamertemplate{enumerate item}
{
  \begin{pgfpicture}{-1ex}{-0.65ex}{1ex}{1ex}
    \usebeamercolor[fg]{item projected}
    {\pgftransformscale{1.75}\pgftext{\normalsize\pgfuseshading{bigsphere}}}
    {\pgftransformshift{\pgfpoint{0pt}{0.5pt}}
      \pgftext{\usebeamercolor[fg]{item projected}\usebeamerfont*{item projected}\insertenumlabel}}
  \end{pgfpicture}%
}

\setbeamertemplate{enumerate subitem}
{
  \begin{pgfpicture}{-1ex}{-0.55ex}{1ex}{1ex}
    \usebeamercolor[fg]{subitem projected}
    {\pgftransformscale{1.4}\pgftext{\normalsize\pgfuseshading{bigsphere}}}
    \pgftext{%
      \usebeamercolor[fg]{subitem projected}%
      \usebeamerfont*{subitem projected}%
      \insertsubenumlabel}
  \end{pgfpicture}%
}

\setbeamertemplate{enumerate subsubitem}
{
  \begin{pgfpicture}{-1ex}{-0.55ex}{1ex}{1ex}
    \usebeamercolor[fg]{subsubitem projected}
    {\pgftransformscale{1.4}\pgftext{\normalsize\pgfuseshading{bigsphere}}}
    \pgftext{%
      \usebeamercolor[fg]{subsubitem projected}%
      \usebeamerfont*{subitem projected}%
      \insertsubsubenumlabel}
  \end{pgfpicture}%
}




\newcommand{\R}{ℝ}
\newcommand{\N}{ℕ}
\newcommand{\Z}{ℤ}
\newcommand{\card}[1]{\lvert{#1}\rvert}
\newcommand{\powerset}[1]{\mathscr{P}(#1)}%\mathscr rather than \mathcal: scr is rounder than cal (at least in XITS Math).
\newcommand{\suchthat}{\;\ifnum\currentgrouptype=16 \middle\fi|\;}
%\newcommand{\Rplus}{\reels^+\xspace}

\AtBeginDocument{%
	\renewcommand{\epsilon}{\varepsilon}
% we want straight form of \phi for mathematics, as recommended in UTR #25: Unicode support for mathematics.
%	\renewcommand{\phi}{\varphi}
}

% with amssymb, but I don’t want to use amssymb just for that.
% \newcommand{\restr}[2]{{#1}_{\restriction #2}}
%\newcommand{\restr}[2]{{#1\upharpoonright}_{#2}}
\newcommand{\restr}[2]{{#1|}_{#2}}%sometimes typed out incorrectly within \set.
%\newcommand{\restr}[2]{{#1}_{\vert #2}}%\vert errors when used within \Set and is typed out incorrectly within \set.
\DeclareMathOperator*{\argmax}{arg\,max}
\DeclareMathOperator*{\argmin}{arg\,min}


%Voting
\newcommand{\allalts}{\mathscr{A}}
\newcommand{\alts}{A}
\newcommand{\feasalts}{F}
\newcommand{\allvoters}{\mathscr{N}}
\newcommand{\voters}{N}
\newcommand{\allsystems}{\mathcal{G}}
\newcommand{\prof}{\mathbf{R}}
\newcommand{\allprofs}{\mathbfcal{R}}
\newcommand{\linors}{\mathcal{L}(\alts)}

\newcommand{\pbasic}[1]{\prof^{#1}_\epsilon}
\newcommand{\pelem}[1]{\prof^{#1}_e}
\newcommand{\pcycl}[1]{\prof^{#1}_c}
\newcommand{\pcycllong}[1]{\prof^{#1}_{cl}}
\newcommand{\pinv}[1]{\overline{\prof_{#1}}}
\newcommand{\dmap}{{\xitsfamily δ}}
%powerset without zero
\newcommand{\powersetz}[1]{\mathcal{P}_\emptyset(#1)}

%logic atom
%⟼ (long)
\DeclareDocumentCommand{\lato}{ O{\prof} O{\alts} }{[#1 \!⟼\! #2]}
%logic atom in
%↝, \stackrel{\in}{\mapsto}, ➲, ⥹
\newcommand{\tightoverset}[2]{%
  \mathop{#2}\limits^{\vbox to -.5ex{\kern-0.9ex\hbox{$#1$}\vss}}}
\DeclareDocumentCommand{\latoin}{ O{\prof} O{\alpha} }{[#1 \tightoverset{\in}{⟼} #2]}
\newcommand{\alllang}{\mathcal{L}}
\newcommand{\ltru}{\texttt{T}}
\newcommand{\lfal}{\texttt{F}}
\newcommand{\laxiom}[1]{{\texgyreherosfamily{\textsc{#1}}}}

%ARG TH
\newcommand{\AF}{\mathcal{AF}}
\newcommand{\labelling}{\mathcal{L}}
\newcommand{\labin}{\textbf{in}\xspace}
\newcommand{\labout}{\textbf{out}}
\newcommand{\labund}{\textbf{undec}\xspace}
\newcommand{\nonemptyor}[2]{\ifthenelse{\equal{#1}{}}{#2}{#1}}
\newcommand{\gextlab}[2][]{
	\labelling{\mathcal{GE}}_{(#2, \nonemptyor{#1}{\ibeatsr{#2}})}
}
\newcommand{\allargs}{A^*}
\newcommand{\args}{A}
\newcommand{\ar}{a}
\newcommand{\ext}{\mathcal{E}}

%MCDA+Arg
\newcommand{\dm}{d}
\newcommand{\ileadsto}{\rightcurvedarrow}
\newcommand{\mleadsto}[1][\eta]{\rightcurvedarrow_{#1}}
\newcommand{\ibeats}{\vartriangleright}
\newcommand{\mbeats}[1][\eta]{\vartriangleright_{#1}}


%MISC
\newcommand{\lequiv}{\Vvdash}
\newcommand{\weightst}{W^{\,t}}

%MCDA classical
\newcommand{\crits}{\mathcal{J}}

%Sorting
\newcommand{\cats}{\mathcal{C}}
\newcommand{\catssubsets}{2^\cats}
\newcommand{\catgg}{\vartriangleright}
\newcommand{\catll}{\vartriangleleft}
\newcommand{\catleq}{\trianglelefteq}
\newcommand{\catgeq}{\trianglerighteq}
\newcommand{\alttoc}[2][x]{(#1 \xrightarrow{} #2)}
\newcommand{\alttocat}[3]{(#2 \xrightarrow{#1} #3)}
\newcommand{\alttoI}{(x \xrightarrow{} \left[\underline{C_x}, \overline{C_x}\right])}
\newcommand{\alttocatdm}[3][t]{\left(#2 \thinspace \raisebox{-3pt}{$\xrightarrow{#1}$}\thinspace #3\right)}
\newcommand{\alttocatatleast}[2]{\left(#1 \thinspace \raisebox{-3pt}{$\xrightarrow[]{≥}$}\thinspace #2\right)}
\newcommand{\alttocatatmost}[2]{\left(#1 \thinspace \raisebox{-3pt}{$\xrightarrow[]{≤}$}\thinspace #2\right)}

\definecolor{darkgreen}{rgb}{0,0.6,0}
\newcommand{\commentOC}[1]{{\selectlanguage{french}{\todo{OC : #1}}}}
%Or: \todo[color=green!40]
\newcommand{\innote}[1]{{\scriptsize{#1}}}

%which line breaks are chosen: accept worse lines, therefore reducing risk of overfull lines. Default = 200
\tolerance=2000
%accept overfull hbox up to...
\hfuzz=2cm
%reduces verbosity about the bad line breaks
\hbadness 5000
%sloppy sets tolerance to 9999
\apptocmd{\sloppy}{\hbadness 10000\relax}{}{}

\bibliographystyle{abbrvnat}
%or \bibliographystyle{apalike} for presentations?

%doi package uses old-style dx.doi url, see 3.8 DOI system Proxy Server technical details, “Users may resolve DOI names that are structured to use the DOI system Proxy Server (http://doi.org (preferred) or http://dx.doi.org).”, https://www.doi.org/doi_handbook/3_Resolution.html
\makeatletter
\patchcmd{\@doi}{dx.doi.org}{doi.org}{}{}
\makeatother

% WRITING
%\newcommand{\ie}{i.e.\@\xspace}%to try
%\newcommand{\eg}{e.g.\@\xspace}
%\newcommand{\etal}{et al.\@\xspace}
\newcommand{\ie}{i.e.\ }
\newcommand{\eg}{e.g.\ }
\newcommand{\mkkOK}{\checkmark}%\color{green}{\checkmark}
\newcommand{\mkkREQ}{\ding{53}}%requires pifont?%\color{green}{\checkmark}
\newcommand{\mkkNO}{}%\text{\color{red}{\textsf{X}}}

\makeatletter
\newcommand{\boldor}[2]{%
	\ifnum\strcmp{\f@series}{bx}=\z@
		#1%
	\else
		#2%
	\fi
}
\newcommand{\textstyleElProm}[1]{\boldor{\MakeUppercase{#1}}{\textsc{#1}}}
\makeatother
\newcommand{\electre}{\textstyleElProm{Électre}\xspace}
\newcommand{\electreIv}{\textstyleElProm{Électre Iv}\xspace}
\newcommand{\electreIV}{\textstyleElProm{Électre IV}\xspace}
\newcommand{\electreIII}{\textstyleElProm{Électre III}\xspace}
\newcommand{\electreTRI}{\textstyleElProm{Électre Tri}\xspace}
% \newcommand{\utadis}{\texorpdfstring{\textstyleElProm{utadis}\xspace}{UTADIS}}
% \newcommand{\utadisI}{\texorpdfstring{\textstyleElProm{utadis i}\xspace}{UTADIS I}}

%TODO
% \newcommand{\textstyleElProm}[1]{{\rmfamily\textsc{#1}}} 


\newlength{\GraphsNodeSep}
\setlength{\GraphsNodeSep}{7mm}

% MCDA Drawing Sorting
\newlength{\MCDSCatHeight}
\setlength{\MCDSCatHeight}{6mm}
\newlength{\MCDSAltHeight}
\setlength{\MCDSAltHeight}{4mm}
%separation between two vertical alts
\newlength{\MCDSAltSep}
\setlength{\MCDSAltSep}{2mm}
\newlength{\MCDSCatWidth}
\setlength{\MCDSCatWidth}{3cm}
\newlength{\MCDSEvalRowHeight}
\setlength{\MCDSEvalRowHeight}{6mm}
\newlength{\MCDSAltsToCatsSep}
\setlength{\MCDSAltsToCatsSep}{1.5cm}
\newcounter{MCDSNbAlts}
\newcounter{MCDSNbCats}
\newlength{\MCDSArrowDownOffset}
\setlength{\MCDSArrowDownOffset}{0mm}

\tikzset{/Graphs/dot/.style={
	shape=circle, fill=black, inner sep=0, minimum size=1mm
}}
\tikzset{/MC/D/S/alt/.style={
	shape=rectangle, draw=black, inner sep=0, minimum height=\MCDSAltHeight, minimum width=2.5cm, anchor=north east
}}
\tikzset{MC/D/S/pref/.style={
	shape=ellipse, draw=gray, thick
}}
\tikzset{/MC/D/S/cat/.style={
	shape=rectangle, draw=black, inner sep=0, minimum height=\MCDSCatHeight, minimum width=\MCDSCatWidth, anchor=north west
}}
\tikzset{/MC/D/S/evals matrix/.style={
	matrix, row sep=-\pgflinewidth, column sep=-\pgflinewidth, nodes={shape=rectangle, draw=black, inner sep=0mm, text depth=0.5ex, text height=1em, minimum height=\MCDSEvalRowHeight, minimum width=12mm}, nodes in empty cells, matrix of nodes, inner sep=0mm, outer sep=0mm, row 1/.style={nodes={draw=none, minimum height=0em, text height=, inner ysep=1mm}}
}}
\tikzset{profile matrix/.style={
	matrix of math nodes, column sep=3mm, row sep=2mm, nodes={inner sep=0.5mm, anchor=base}
}}

% GUI
\tikzset{/GUI/button/.style={
	rectangle, very thick, rounded corners, draw=black, fill=black!40%, top color=black!70, bottom color=white
}}

% Beliefs
\tikzset{/Beliefs/D/S/attacker/.style={
	shape=rectangle, draw, minimum size=8mm
}}
\tikzset{/Beliefs/D/S/supporter/.style={
	shape=circle, draw
}}

\newcommand{\tikzmark}[1]{%
	\tikz[overlay, remember picture, baseline=(#1.base)] \node (#1) {};%
}


\begin{acronym}
\acro{AMCD}{Aide Multicritère à la Décision}
\acro{ASA}{Argument Strength Assessment}
\acro{DA}{Decision Analysis}
\acro{DM}{Decision Maker}
\acro{DPr}{Deliberated Preferences}
\acro{DRSA}{Dominance-based Rough Set Approach}
\acro{DSS}{Decision Support Systems}
\acrodefplural{DSS}{Decision Support Systems}
% \newacroplural{DSS}[DSSes]{Decision Support Systems}
\acro{EJOR}{European Journal of Operational Research}
\acro{LNCS}{Lecture Notes in Computer Science}
\acro{MCDA}{Multicriteria Decision Aid}
\acro{MIP}{Mixed Integer Program}
\acro{NCSM}{Non Compensatory Sorting Model}
\acro{PL}{Programme Linéaire}
\acro{PLNE}{Programme Linéaire en Nombres Entiers}
\acro{PM}{Programme Mathématique}
\acro{MP}{Mathematical Program}
\acro{MIP}{Mixed Integer Program}
% \newacroplural{PM}{Programmes Mathématiques}
%acrodefplural since version 1.35, my debian has \ProvidesPackage{acronym}[2009/01/25, v1.34, Support for acronyms (Tobias Oetiker)]
\acrodefplural{PM}{Programmes Mathématiques}
\acro{PMML}{Predictive Model Markup Language}
\acro{RESS}{Reliability Engineering \& System Safety}
\acro{SMAA}{Stochastic Multicriteria Acceptability Analysis}
\acro{URPDM}{Uncertainty and Robustness in Planning and Decision Making}
\acro{XML}{Extensible Markup Language}
\end{acronym}


\setbeamertemplate{headline}[singleline]

\title{Axiomatics, then what?}
\subject{Social Choice}
\keywords{Axiomatization}
\author{Olivier Cailloux}
\institute[LAMSADE]{LAMSADE, Université Paris-Dauphine}
\date{\formatdate{30}{1}{2017}}

\begin{document}
\begin{frame}[plain]
	\tikz[remember picture,overlay]{
		\path (current page.south west) node[anchor=south west, inner sep=0] {
			\includegraphics[height=1cm]{LAMSADE95.jpg}
		};
		\path (current page.south) ++ (0, 1mm) node[anchor=south, inner sep=0] {
			\includegraphics[height=9mm]{Dauphine.jpg}
		};
		\path (current page.south east) node[anchor=south east, inner sep=0] {
			\includegraphics[height=1cm]{PSL.png}
		};
	}
	\titlepage
\end{frame}
\addtocounter{framenumber}{-1}

\section{Axiomatics}
\subsection{Axiomatics}
\begin{frame}
	\frametitle{\subsecname}
	\begin{quote}
		Rather than dream up a multitude of arbitration schemes and determine whether or not each withstands the best of plausibility in a host of special cases, let us invert the procedure. Let us examine our subjective intuition of fairness and formulate this as a set of precise desiderata that any acceptable arbitration scheme must fulfil. Once these desiderata are formalized as axioms, then the problem is reduced to a mathematical investigation of the existence of and characterization of arbitration schemes which satisfy the axioms.
	\end{quote}
	\citet[p. 121]{luce_games_1957}
\end{frame}

\subsection{Beyond axioms?}
\begin{frame}
	\frametitle{\subsecname}
	Possible view (that I want to counter):
	\begin{itemize}
		\item Mathematically capture the behavior of a rule (with axioms…)
		\item Left out for the user: confront their intuition about fairness with the axioms
		\item Scientific approach stops at the first step
		\item The rest is ultimately subjective
	\end{itemize}
	I propose a different view: 
	\begin{itemize}
		\item Confronting our subjective intuition of fairness with axioms is hard
		\item Nobody can do this
		\item Because we can’t compute implications
		\item We can help doing this scientifically
	\end{itemize}
\end{frame}

\section{The hard problem of reaching the equilibrium}
\subsection{Problem reduced?}
\begin{frame}
	\frametitle{\subsecname}
	\begin{itemize}
		\item Formulate our subjective intuition of fairness as axioms?
		\item \emph{Excluding} impossibilities?
		\item Making sure no hidden axiom is left out?
		\item Accept Arrow’s axioms?
		\item Who’s subjective intuition?
	\end{itemize}
\end{frame}

\subsection{Understand the rule?}
\begin{frame}
	\frametitle{\subsecname}
	\begin{itemize}
		\item When do you \emph{understand} a voting rule?
		\item Borda rule: I know how to count scores
		\item Is that all?
		\item Recall we want to capture our (?) idea of fairness
	\end{itemize}
\end{frame}

\subsection{What do axioms say?}
\begin{frame}
	\frametitle{\subsecname}
	\begin{itemize}
		\item Humans have limited deductive power
		\item Hence, knowing the definition of a rule does not determine whether we accept it
		\item Knowing the \emph{axiomatization} of a rule does not determine whether we accept it
		\item Wanted: equilibrium between principles and case-based, concrete intuitions \citep{goodman_fact_1983, rawls_theory_1999}
	\end{itemize}
\end{frame}

\subsection{What do axioms say? (Illustrations)}
\begin{frame}
	\frametitle{\subsecname}
	\begin{example}[Borda]
		\begin{itemize}
			\item We know Borda is the rule that satisfies neutrality, reinforcement, faithfulness, cancellation
			\item This does not obviously say that Borda fails on Condorcet
		\end{itemize}
	\end{example}
	\begin{example}[Dictatorship]
		\begin{itemize}
			\item One may want Arrow’s axioms
			\item But fail to see what it implies: dictatorship
			\item Once the implications are understood, one does not want all of Arrow’s axioms any more
		\end{itemize}
	\end{example}
\end{frame}

\subsection{Fishburn-against-Condorcet}
\begin{frame}[fragile]
	\frametitle{\subsecname}
	\begin{minipage}{5.5cm}
		\citet[p. 544]{fishburn_paradoxes_1974} argument against the Condorcet principle (see also \url{http://rangevoting.org/FishburnAntiC.html}).
		\begin{block}{Condorcet winner}
%		\setlength\abovedisplayskip{0 ex}
%		\begin{equation}
			$w \text{ VS } \mu, \mu \in \set{a, \ldots, h}\text{? }\uncover<2->{51/101}$
%		\end{equation}
		\end{block}
	\end{minipage}%
	\begin{minipage}{\columnwidth-5cm}
		\small
		\begin{equation}
			\begin{array}{lrrrrrr}
				&\multicolumn{6}{c}{\text{nb voters}}\\
			\cmidrule{2-7}
					&31	&19	&10	&10	&10	&21	\\
			\midrule
				1	&a	&a	&f	&g	&h	&h	\\
				2	&b	&b	&w	&w	&w	&g	\\
				3	&c	&c	&a	&a	&a	&f	\\
				4	&d	&d	&h	&h	&f	&w	\\
				5	&e	&e	&g	&f	&g	&a	\\
				6	&w	&f	&e	&e	&e	&e	\\
				7	&g	&g	&d	&d	&d	&d	\\
				8	&h	&h	&c	&c	&c	&c	\\
				9	&f	&w	&b	&b	&b	&b	\\
			\end{array}
		\end{equation}
	%	\begin{tikzpicture}
	%		\path node[profile matrix] (profile) {
	%				&31	&19	&10	&10	&10	&21	\\
	%			1	&a	&a	&f	&g	&h	&h	\\
	%			2	&b	&b	&w	&w	&w	&g	\\
	%			3	&c	&c	&a	&a	&a	&f	\\
	%			4	&d	&d	&h	&h	&f	&w	\\
	%			5	&e	&e	&g	&f	&g	&a	\\
	%			6	&w	&f	&e	&e	&e	&e	\\
	%			7	&g	&g	&d	&d	&d	&d	\\
	%			8	&h	&h	&c	&c	&c	&c	\\
	%			9	&f	&w	&b	&b	&b	&b	\\
	%		};
	%	\end{tikzpicture}
	%	BUT ranks:
	\end{minipage}
	\vspace{-1pt}
	\onslide<3>
	\begin{equation}
		\vspace*{-8.4pt}
		\begin{array}{lrrrrrrrrr}
			&\multicolumn{9}{c}{\text{ranks}}\\
		\cmidrule{2-10}
			&1	&≤2	&≤3	&≤4	&≤5	&≤6	&≤7	&≤8	&≤9	\\
		\midrule
		w	&0	& 30	& 30	& 51	& 51	& 82	& 82	& 82	&101	\\
		a	&50	& 50	& 80	& 80	& 101	& 101	& 101	&101	&101	\\
		\end{array}
	\end{equation}
\end{frame}

\section{Possible approach}
\subsection{Possible approach}
\begin{frame}
	\frametitle{\subsecname}
	\begin{itemize}
		\item To know whether one accepts a voting rule, we have to check whether one accepts the implications of the voting rule
		\item Can’t be done exhaustively
		\item Can be done using known possible problematic cases
		\item Axiomatics can help in providing reasonings
		\item Different axiomatics may have different convincing power
		\item Acceptance of people in concrete cases can in principle be studied empirically \citep{gaertner_empirical_2012}
	\end{itemize}
	\begin{block}{Proposition}
		\begin{itemize}
			\item Proposition: propose a research program aiming at such a study / or
			\item Propose first steps (a framework?)
		\end{itemize}
	\end{block}
\end{frame}

\subsection{Key element}
\begin{frame}
	\frametitle{\subsecname}
	\begin{itemize}
		\item Obviously, in many cases, no single decisive rule is the most appropriate
		\item We have to allow for incompleteness
		\item E.g. Impossible to completely rank all universities by “quality”
		\item Possible approach: search for non decisive rules (possible winners…)
	\end{itemize}
\end{frame}

\begin{frame}[plain]
	\addtocounter{framenumber}{-1}
	\begin{center}
		\huge
		\textit{Thank you for your attention!}
	\end{center}
\end{frame}

\appendix
\AtBeginSection{
}

\section{Bibliography}
\begin{frame}[allowframebreaks]
	\frametitle{\secname}
	\def\newblock{\hskip .11em plus .33em minus .07em}
 	\bibliography{zotero,philo-eco,beliefs}
\end{frame}

\section{Voting rules}
\subsection{Definition}
\begin{frame}
	\frametitle{\subsecname}
	\begin{itemize}
		\item Voting rule: a systematic way of aggregating different opinions and decide
		\item Multiple reasonable ways of doing this
		\item Different voting rules have different interesting properties
		\item None satisfy all desirable properties
	\end{itemize}
\end{frame}

\begin{frame}[fragile]
	\frametitle{Voting rule}
	
	\begin{description}[Alternatives]
		\item[Alternatives] $\allalts = \set{a, b, c, d, \ldots}$
		\item[Voters] $\allvoters = \set{1, 2, \ldots}$
		\item[Profile] function $\prof$ from $\allvoters$ to linear orders on $\allalts$.
		\item[Voting rule] function $f$ mapping each $\prof$ to winners $\emptyset \subset \alts \subseteq \allalts$.
	\end{description}
	\vfill
	\begin{center}
		\begin{tikzpicture}
			\path node[profile matrix] (profile) {
				\prof(1)&
				\prof(2)
				\\
				| (profile11) | a&
				b
				\\
				b&
				a
				\\
				c&
				| (profile32) | c
				\\
			};
			\path ($(profile.south west)!.5!(profile.south east)$) ++ (0, -5mm) node {$\prof$};
			
			\path node[draw, rectangle, fit=(profile11) (profile32), outer xsep=2mm, outer ysep=1mm] (justprofile) {};
			\path (justprofile.east) ++ (2.5cm, 0) node[inner sep=0] (winners) {\mbox{} $\alts = \Set{a, b}$};
			\path[draw, ->] (justprofile.east) to[bend left=35] node[anchor=south] {$f$} (winners.west);

			\path[draw, decorate, decoration={brace, mirror}] (justprofile.south west) -- (justprofile.south east);
		\end{tikzpicture}
	\end{center}
\end{frame}

\subsection{Example of a profile}
\begin{frame}[fragile]
	\frametitle{\subsecname}
	\begin{center}
		$\begin{array}{lrrrrrr}
			&\multicolumn{6}{c}{\text{nb voters}}\\
		\cmidrule{2-7}
				&33	&16	&3	&8	&18	&22	\\
		\midrule
			1	&a	&b	&c	&c	&d	&e	\\
			2	&b	&d	&d	&e	&e	&c	\\
			3	&c	&c	&b	&b	&c	&b	\\
			4	&d	&e	&a	&d	&b	&d	\\
			5	&e	&a	&e	&a	&a	&a	\\
		\end{array}$
	\end{center}
	Who wins?
	\tikz[remember picture,overlay]{
		\path (current page.south east) ++ (0, 1em) node[anchor=south east, inner sep=0] {
			\tiny Thanks to \href{http://www.lamsade.dauphine.fr/~lang/}{Jérôme Lang}
		};
	}
	\begin{itemize}
		\item Most top-1: $a$
		\item $c$ is in the top 3 for everybody
		\item delete worst first, lowest nb of pref: $c$, $b$, $e$, $a$ ⇒ $d$
		\item delete worst first, from bottom: $a$, $e$, $d$, $b$ ⇒ $c$
		\item Borda: $b$
		\item Condorcet: $c$
	\end{itemize}
\end{frame}

\subsection{Two voting rules}
\begin{frame}
	\frametitle{Borda}
	
Given a profile $\prof$:
	\begin{itemize}
		\item score of $a \in \allalts$: number of alternatives it beats
		\item the highest scores win
	\end{itemize}
	
	\begin{equation}
		\prof =
		\begin{array}{rrrrr}
			a	&	a	&	a	&	b	&	b\\
			b	&	b	&	b	&	c	&	c\\
			c	&	c	&	c	&	a	&	a
		\end{array}
	\end{equation}
	\begin{itemize}
		\item score $a$ is\dots? \pause $2 + 2 + 2 = 6$
		\item score $b$ is $1 + 1 + 1 + 2 + 2 = 7$
		\item score $c$ is $1 + 1 = 2$
	\end{itemize}
	Winner: $b$.
\end{frame}

\begin{frame}
	\frametitle{Condorcet’s principle}
	\begin{block}{Condorcet’s principle}
		We ought to take the Condorcet winner as sole winner if it exists.
		\begin{itemize}
			\item $a$ \emph{beats} $b$ iff more than half the voters prefer $a$ to $b$.
			\item $a$ is a \emph{Condorcet winner} iff $a$ beats every other alternatives.
		\end{itemize}
	\end{block}
	\vfill
	\begin{equation}
		\prof =
		\begin{array}{rrrrr}
			a	&	a	&	a	&	b	&	b\\
			b	&	b	&	b	&	c	&	c\\
			c	&	c	&	c	&	a	&	a
		\end{array}
	\end{equation}
	 Who wins? \pause $a$
\end{frame}

\begin{frame}
	\frametitle{Condorcet’s principle and a voting rule}
	\begin{itemize}
		\item Condorcet’s principle does not define a voting rule. Why? \pause
		\item No winner is defined when no Condorcet winner
	\end{itemize}
	\begin{equation}
		\prof =
		\begin{array}{rrr}
			b	&	c	&	d\\
			c	&	d	&	b\\
			a	&	b	&	a\\
			d	&	a	&	c
		\end{array}
	\end{equation}
	\pause
	$a$ loses against $b$; $b$ against $d$; $c$ against $b$; $d$ against $c$
	\begin{itemize}
		\item Dodgson’s method (1876): candidates “closest” to being Condorcet winners {\tiny (in nb of swaps)}
	\end{itemize}
\end{frame}

\end{document}

\section{Frame template}
\subsection{}
\begin{frame}
	\frametitle{\subsecname}
	\begin{itemize}
		\item 
	\end{itemize}
\end{frame}

